\section{Introduction}
\label{sec:introduction}

Online large-language-model (LLM) serving must control both time to first token
(TTFT) and time per output token (TPOT). Prefill processes the prompt and is
typically compute intensive; decode repeatedly reads a growing KV cache and is
often memory intensive. A scheduler therefore needs to advance queued prompts
without creating long gaps for active generations.

Continuous batching rebuilds the batch at every iteration~\cite{yu2022orca,
kwon2023vllm}. Chunked prefill further limits prompt work and mixes a bounded
prefill chunk with decode tokens, reducing generation stalls~\cite{
agrawal2024sarathi,holmes2024fastgen}. This bound controls how much prefill work
enters one iteration, but does not independently specify how often prefill is
served: under a decode-first token budget, available prefill capacity depends
on the current decode batch.

Physical prefill--decode disaggregation can provision the phases independently,
but requires distinct instances and KV transfer~\cite{patel2024splitwise,
zhong2024distserve}. Spatial multiplexing instead depends on concurrent kernels
and controllable compute partitioning. We study the complementary deployment
budget in which one tensor-parallel (TP) replica occupies the allocated devices
and both phases retain its weights and KV cache locally.

We present \system{}, a scheduler that separates state sharing from execution.
It preserves chunked continuation, but each model iteration is prefill-only or
decode-only. A selector exposes the long-run prefill service tendency as a
configurable share, and work-conserving fallback serves the nonempty phase.
This creates an explicit interface between TTFT and TPOT pressure. The benefit
is conditional: prefill-only slots can reduce queueing when decode has TPOT
slack, whereas they can harm a saturated decode queue.

Our contributions are: (1) identifying the chunk bound and cross-iteration
phase share as distinct controls; (2) implementing bounded, phase-pure
scheduling in a shared TP replica without weight duplication or KV transfer;
and (3) evaluating both its gains and boundary on two Ascend 910B3 NPUs.
Three-seed Azure-derived bursts show strict joint-SLO improvements, a large
reconstructed MaaS workload supports the TTFT trend, and decode-heavy cases
show that PD-TDM is not a universal replacement for mixed chunked prefill.
